\documentclass[12pt]{article}
\newcommand{\bottomMargin}{2cm}

\newcommand{\header}{
	ДИЗАЙН И АНАЛИЗ НА АЛГОРИТМИ - \\ ПРАКТИКУМ\\
	\vspace{0.1cm}
	Летен семестър, 2026 г., второ контролно\\
	\vspace{0.1cm}
}
\newcommand{\tl}{$0.4$ сек.}
\newcommand{\ml}{$256$ MB}

\input{structure.tex}
\usepackage{newunicodechar}
\newunicodechar{ѝ}{{$\grave{\text{и}}$}}
\newunicodechar{Ѝ}{{$\grave{\text{И}}$}}
\begin{document}
	
	\section{Задача K5. Bangaranga}
    
    Новият хит на Дара „Bangaranga“ взривява класациите, а заснемането на официалното видео изисква перфектна и сложна хореография върху специална интерактивна сцена.

    Сцената е конструирана като квадратна мрежа от светодиодни (LED) панели с размер $N \times N$. Между самите панели има разделителни линии - $N-1$ в посока север-юг и $N-1$ в посока изток-запад. За да премине от един панел в съседен (на север, изток, юг или запад) в пълен синхрон с бързия електронен бийт, на Дара ѝ отнема точно $T$ единици време.
    
    В началото на видеото Дара стартира своето движение от панела в горния ляв ъгъл, а финалната поза на хореографията трябва да бъде заета точно пред главната камера в долния десен ъгъл. Тъй като песента има силно изявени вокални акценти, на всяко трето преместване Дара трябва да спре на текущия панел (без да се брои началната ѝ позиция, но може да се включи финалният панел), за да изпее емблематичната реплика: „bangaranga“. Всеки панел има различна интерактивна графика и светлинни ефекти, които изискват специфично време за изпълнение на вокалната поза, затова времето за спиране варира в зависимост от панела.
    
    Моля, помогнете на режисьора на клипа, като напишете програма \textbf{bangaranga}, която да определи минималното време, необходимо на Дара да стигне до финалната позиция пред камерата, изпълнявайки точно вокалните пози на всеки три стъпки.
    
    \subsection*{Вход}
    На първия ред на стандартния вход се въвеждат $N$ и $T$. 
    Следващите $N$ реда съдържат по $N$ цели положителни числа, описващи времето, необходимо за вокалната поза на всеки LED панел. Първото число на първия ред съответства на северозападния ъгъл.
    
    \subsection*{Изход}
    На стандартния изход изведете едно число – минималното време, необходимо на Дара да завърши хореографията в долния десен ъгъл.
    
    \subsection*{Ограничения}
    \begin{itemize}
        \item $3 \leq N \leq 500$
        \item $0 \leq T \leq 1~000~000$
        \item Времето за изпълнение на вокалната поза на всеки панел не надвишава $100~000$.
    \end{itemize}

    \subsection*{Примери}
    \begin{table}[ht]
        \begin{tblr}{|X[30,l]|X[8,l]|X[62,l]|}
            \hline
            \textbf{Вход} & \textbf{Изход} & \textbf{Обяснение на примера} \\
            \hline
            \texttt{4 2\\30 92 36 10\\38 85 60 16\\41 13 5 68\\20 97 13 80} & 
            \texttt{31} &
            {Оптималното решение за този пример включва придвижване надясно с 3 панела (където Дара прави първото си спиране за вокална поза с време $10$), след това преместване надолу 2 пъти и наляво веднъж (където прави второто си спиране с време $5$), и накрая придвижване надолу и надясно до финалната камера.}\\
            \hline
        \end{tblr}
    \end{table}
	
\end{document}
